Mit der grossen Achse nimmt auch die grosse Halbachse um \pgaO\ zu: $a_2 = \left(1+\pgaO\right)a_1$. Nach dem dritten Kepler'schen Gesetz ist $T^2 \propto a^3$, also
\begin{align}
 \frac{T_2}{T_1} &= \qTaF = \left(1+\pgaO\right)^{3/2} = \qTa
\end{align}
Die Umlaufzeit \result{\text{nimmt um } \pdTP{0}{2} \text{ zu}}.
